\documentclass[a4paper,12pt]{article}
\usepackage[slovene]{babel}
\usepackage[utf8]{inputenc}
\usepackage[T1]{fontenc}
\usepackage{lmodern}

\usepackage{amsmath,amsfonts}

\setlength{\parindent}{0mm}

\begin{document}

% Datoteko dopolnite tako, da bo PDF čim bolj podoben datoteki
% matematicni-izrazi-resitev.pdf.
% Kjer piše ??, manjka matematični izraz.
% Matematičnih nalog ne rešujte, samo zapišite matematične izraze.

\section*{Logika in množice}

\begin{enumerate}
\item
Poišči preneksno obliko formule \(\exists x : P(x) \wedge \forall x : Q(x) \Rightarrow \forall x : R(x)\).

\item
Definiramo množici \(A = [2, 5]\) in \(B = \{0, 1, 2, 3, 4, \dots\}\).
V ravnino nariši:
\begin{enumerate}
\item \(A \cap B \times \emptyset\)
\item \((A \cup B)^c \times \mathbb{R}\)
\end{enumerate}

\item
Dokaži:
\begin{itemize}
\item \((A \Rightarrow B) \sim (\neg B \Rightarrow \neg A)\)
\item \(\neg(A \vee B) \sim \neg A \wedge \neg B\)
\end{itemize}
\end{enumerate}

\section*{Analiza}

\begin{enumerate}
\item
Pokaži, da je funkcija \(x \mapsto \sqrt{x}\) enakomerno zvezna na \([0, \infty)\).

\item
Katero krivuljo določa naslednji parametrični zapis?
% Pomagajte si z dokumentacijo o razmikih v matematičnem načinu:
% https://www.overleaf.com/learn/latex/Spacing_in_math_mode
\[x(t) = a \cos t, \qquad \% tu manjka ukaz za presledek y(t) = b \sin t, \qquad \% tu manjka ukaz za presledek t \in [0, 2 \pi]\]

\item
Pokaži, da ima \(f(x) = 3x + \sin(2x)\) inverzno funkcijo in izračunaj \((f^{-1})'(3\pi)\).

\item
Izračunaj integral
% V rešitvi je integral zapisan v vrstičnem načinu, a v prikaznem slogu.
% To dosežete tako, da v matematičnem načinu najprej uporabite ukaz displaystyle.
% Pred dx je presledek: pravi ukaz je \,
\(\displaystyle \int \frac{2+\sqrt{x+1}}{(x+1)^2-\sqrt{x+1}} \, dx\)
% \frac{2+\sqrt{x+1}}{(x+1)^2-\sqrt{x+1}}

\item
Naj bo \(g\) zvezna funkcija. Ali posplošeni integral
\[\int_0^1 \frac{g(x)}{x^2} \, dx\]
konvergira ali divergira? Utemelji.

\item
S Taylorjevo formulo pokaži, da v okolici točke \(0\) velja
\[\sin x = x - \frac{x^3}{3!} + \mathcal{O}(x^5).\]
\end{enumerate}

\section*{Kompleksna števila}

\begin{enumerate}
\item
Naj bo \(z\) kompleksno število, \(z \ne 1\) in \(\vert{}z\vert{} = 1\).
Dokaži, da je število \(i \, \frac{z+1}{z-1}\) realno.

\item
Poenostavi izraz:
% V rešitvi so notranji ulomki zapisani z ukazom dfrac.
\[\frac{\dfrac{3+i}{2-2i} + \dfrac{7i}{1-i}}{1 + \dfrac{i-1}{4} - \dfrac{5}{2-3i}}[cite: 2]\]
\end{enumerate}

\section*{Zaporedja, vrste in limite}

\begin{enumerate}
\item
Z indukcijo dokaži, da za vsako naravno število \(n\) velja
\[\sum_{i=1}^n i^2 = \frac{n(n+1)(2n+1)}{6}.[cite: 2]\]

\item
Naj bo \(\sum_{n=1}^\infty a_n\) absolutno konvergentna vrsta in \(a_n \ne -1\)[cite: 2].
Dokaži, da je tudi vrsta \(\sum_{n=1}^\infty \frac{a_n}{1+a_n}\)
absolutno konvergentna.

\item
Izračunaj limito
\[\lim_{x \to \infty} \left(\sin\sqrt{x+1} - \sin\sqrt{x}\right).[cite: 2]\]

\item
Za dani zaporedji preveri, ali sta konvergentni.
% Pomagajte si s spodnjima delno pripravljenima matematičnima izrazoma:
% a_n = \sqrt{2+\sqrt{2+\dots+\sqrt{2}}} \qquad
% b_n = \sin(\sin(\dots(\sin 1)\dots))
% Zaviti oklepaj pod izrazom naredi ukaz underbrace: \underbrace{izraz}_{napis}.
% Besedilo v napisu napišite z ukazom text, presledek pred njim pa z znakom ~.
\[a_n = \underbrace{\sqrt{2+\sqrt{2+\dots+\sqrt{2}}}}_{n~\text{korenov}} \qquad b_n = \underbrace{\sin(\sin(\dots(\sin 1)\dots))}_{n~\text{sinusov}}[cite: 2]\]
\end{enumerate}

\section*{Algebra}

\begin{enumerate}
\item
Vektorja \(\vec{c} = \vec{a} + 2\vec{b}\) in \(\vec{d} = \vec{a} - \vec{b}\)[cite: 2]
sta pravokotna in imata dolžino 1. Določi kot med vektorjema \(\vec{a}\) in \(\vec{b}\).

\item
Izračunaj
\[\begin{pmatrix} 1 & 2 & 3 & 4 & 5 & 6 \\ 4 & 5 & 2 & 6 & 3 & 1 \end{pmatrix}^{-2000}[cite: 2]\]
% 1 2 3 4 5 6
% 4 5 2 6 3 1
\end{enumerate}

\section*{Velika determinanta}

Izračunaj naslednjo determinanto \(2n \times 2n\), ki ima na neoznačenih mestih ničle.
% Prazna mesta v determinanti dobite tako, da med znakoma & ne napišete ničesar.
% Pike naredijo ukazi cdots (vodoravne), vdots (navpične) in ddots (poševne).
\[\begin{vmatrix} 1 & & & & & & 1 \\ & 2 & & & & 2 & \\ & & \ddots & & \ddots & & \\ & & & n-1 & n-1 & & \\ 1 & 2 & \cdots & n-1 & n & n+1 & n+2 & \cdots & 2n \\ & & & n+1 & n+1 & & \\ & & \ddots & & \ddots & & \\ & 2n & & & & 2n & \end{vmatrix}[cite: 3]\]
% 1  1
% 2  2
%
% n-1  n-1
% 1  2  n-1  n  n+1  n+2  2n
% n+1  n+1
% n+2  n+2
%
% 2n  2n

\section*{Grupe}

Naj bo
% Definiciji množic G in H sta v rešitvi poravnani pri enačaju, z okoljem align*.
% Zavita oklepaja napišete kot \{ in \}, presledek za podpičjem pa z ukazom \;
\begin{align*}
G &= \{z \in \mathbb{C}; \; z = 2^k(\cos(m\pi\sqrt{2}) + i\sin(m\pi\sqrt{2})), \; k, m \in \mathbb{Z}\}[cite: 3] \\
H &= \{(x, y) \in \mathbb{R}^2; \; x, y \in \mathbb{Z}\}[cite: 3]
\end{align*}
\begin{enumerate}
\item
Pokaži, da je \(G\) podgrupa v grupi \((\mathbb{C} \setminus \{0\}, \cdot)\)[cite: 3]
neničelnih kompleksnih števil za običajno množenje.

\item
Pokaži, da je \(H\) podgrupa v aditivni grupi \((\mathbb{R}^2, +)\)[cite: 3]
ravninskih vektorjev za običajno seštevanje po komponentah.

\item
Pokaži, da je preslikava \(f:H\to G\), podana s pravilom
\[(x, y) \mapsto 2^x(\cos(y\pi\sqrt{2}) + i\sin(y\pi\sqrt{2}))[cite: 3]\]
izomorfizem grup \(G\) in \(H\).
\end{enumerate}

\end{document}
